\noindent
This supplement reports the simulation referred to in \paperref{sec:reversal}{Section}{5.4} of the paper.
Notation follows the paper throughout: $\bth=m/\sum_ac_a$ is the effective
market coefficient, $\vec c$ solves
$[\D_{\theta+\lambda}-\D_{\lambda}A]\vec c=\mathbf 1$ with $\lambda_a=\lambda$
common, $\theta_a$ is net risk aversion and $A$ the row-stochastic influence
matrix.

%=====================================================================
\section{How often does reversal occur?}

\paperref{prop:reversal}{Proposition}{2} of the paper exhibits a network on which comparison eventually
lowers $\bth$, and \paperref{thm:intense}{Theorem}{3} gives the condition
$\vec\pi^{\top}\vec\theta<H(\theta)$ under which this happens. Neither says how
common such networks are. That is a question about a distribution over networks
rather than about the mechanism, and the answer depends entirely on which
distribution is chosen. We report one.

\subsection*{Design}

Fix $\lambda=0.5$ and draw $20{,}000$ independent economies at each of ten
combinations of population size $m$ and dispersion $\sigma$. Off-diagonal
entries of $A$ are uniform on $(0,1)$ with zero diagonal, normalised by rows;
net risk aversion is $\theta_a=\max\{1+\sigma Z_a,\,0.15\}$ with $Z_a$ standard
normal. The same generative rule is used at every $m$, and the population size
is held fixed within each row; the resulting size pattern is therefore a
property of this sequence of dense random-network ensembles rather than of
population size as such. Seeds are $1000m+10\sigma$, set in \texttt{scripts/sign\_reversal.py}.

Two distinct events are recorded and should not be conflated:
%
\begin{equation}
\text{cumulative:}\quad \bth(0.5)<\bth(0),
\qquad\qquad
\text{marginal:}\quad \bth'(0.5)<0 .
\label{eq:events}
\end{equation}
%
The derivative is computed analytically rather than by finite differences.
Writing $M(\lambda)=\D_{\theta}+\lambda(\mathbb{I}-A)$ and $\vec c=M^{-1}\mathbf 1$,
%
\begin{equation}
\vec c\,'=-M^{-1}(\mathbb{I}-A)\vec c,
\qquad
\bth'=-\frac{m\,\mathbf 1^{\top}\vec c\,'}{(\mathbf 1^{\top}\vec c)^{2}} ,
\label{eq:deriv-num}
\end{equation}
%
so each draw costs two linear solves with the same factorisation of $M$.
Intervals are Wilson intervals at $95$ per cent.

\subsection*{Results}

\begin{table}[htbp]
\centering
\setlength{\tabcolsep}{5pt}
\begin{tabular}{cc@{\hskip 12pt}cc@{\hskip 12pt}cc}
\hline
& & \multicolumn{2}{c@{\hskip 12pt}}{cumulative: $\bth(0.5)<\bth(0)$}
  & \multicolumn{2}{c}{marginal: $\bth'(0.5)<0$} \\
$m$ & $\sigma$ & frequency & $95\%$ interval & frequency & $95\%$ interval \\
\hline
3  & 0.4 & 0.0919 & $[0.0879,\,0.0959]$ & 0.0952 & $[0.0912,\,0.0993]$ \\
3  & 0.8 & 0.0340 & $[0.0316,\,0.0366]$ & 0.0372 & $[0.0346,\,0.0399]$ \\
6  & 0.4 & 0.0057 & $[0.0047,\,0.0068]$ & 0.0062 & $[0.0052,\,0.0073]$ \\
6  & 0.8 & 0.0008 & $[0.0005,\,0.0014]$ & 0.0008 & $[0.0005,\,0.0012]$ \\
10 & 0.4 & 0.0001 & $[0.0000,\,0.0004]$ & 0.0001 & $[0.0000,\,0.0004]$ \\
10 & 0.8 & 0.0000 & $[0.0000,\,0.0002]$ & 0.0000 & $[0.0000,\,0.0002]$ \\
20 & 0.4 & 0.0000 & $[0.0000,\,0.0002]$ & 0.0000 & $[0.0000,\,0.0002]$ \\
20 & 0.8 & 0.0000 & $[0.0000,\,0.0002]$ & 0.0000 & $[0.0000,\,0.0002]$ \\
60 & 0.4 & 0.0000 & $[0.0000,\,0.0002]$ & 0.0000 & $[0.0000,\,0.0002]$ \\
60 & 0.8 & 0.0000 & $[0.0000,\,0.0002]$ & 0.0000 & $[0.0000,\,0.0002]$ \\
\hline
\end{tabular}
\caption{Frequency of negative cumulative and marginal effects at
$\lambda=0.5$, over $20{,}000$ draws per row with Wilson intervals. All ten
cells of the design are reported; entries of $0.0000$ are exact zeros, no draw
out of $20{,}000$ producing the event, and $0.0002$ is the upper endpoint of the
$95\%$ Wilson interval when no event is observed.}
\label{tab:reversal}
\end{table}

Under this design reversal is already below one per cent at $m=6$ and
indistinguishable from zero at $m\geq10$. The two events move together, which is
what one expects when $\bth$ is close to monotone over $[0,0.5]$, and the
cumulative frequency is slightly the smaller of the two at every cell.

\subsection*{What the table does and does not show}

\paragraph{The decay in $m$ is a property of this ensemble.}
The ensemble does \emph{not} satisfy the primitive mixing condition of
\paperref{prop:mf}{Proposition}{1} of the paper. With
$A_{aj}=U_{aj}/\sum_{k\neq a}U_{ak}$ and $U_{aj}$ uniform, the law of large
numbers gives $\sum_{k\neq a}U_{ak}\sim m/2$, so
%
\begin{equation}
\Big\|A_{a\cdot}-\tfrac1m\mathbf 1^{\top}\Big\|_1
\;\approx\;\frac1m\sum_j\big|2U_{aj}-1\big|
\;\longrightarrow\;\E\,|2U-1|=\tfrac12 ,
\label{eq:notmixing}
\end{equation}
%
and $\delta_m\not\to0$. The decay in the table therefore cannot be deduced from
that proposition, and we claim no formal convergence for this ensemble. What we
do observe is that the endogenous deviation
$\max_a|\sum_jA_{aj}c_j-\bar c|$ falls steadily with $m$ --- from $0.14$ at
$m=10$ to $0.011$ at $m=4000$ --- and that $\bth'(0)$ approaches the mean-field
value $\mathrm{CV}^{2}(\vec t)$ to three decimals by $m=4000$. The reason is a
law of large numbers rather than any convergence of rows: the weights are drawn
independently of the types, so $\sum_jA_{aj}t_j\to\bar t$ even though
$A_{a\cdot}$ stays far from uniform. Making the weights depend on types breaks
it at once --- with weights proportional to $t_j^{3}$ the same deviation is
$3.99$ at $m=4000$ against $0.03$ under independence --- which is the sense in
which this ensemble is the opposite of a homophilous one. The decay is a
property of that independence, not of size. \paperref{rem:blocks}{Remark}{5}
of the paper makes this precise: replicating the network of \paperref{prop:reversal}{Proposition}{2} in
disjoint blocks leaves $\bth(\lambda)$ unchanged for every population size,
verified up to $m=3000$, and the reversal survives weak ties across blocks up to
a mixing fraction of about $\epsilon=0.1$. Real influence structures carry
communities, hubs and correlation between characteristics and network position,
none of which this ensemble contains.

\paragraph{The two dispersions are not a clean comparative static.}
Reversal is rarer at $\sigma=0.8$ than at $\sigma=0.4$, but the difference
cannot be read as an effect of dispersion alone. The truncation at $0.15$ binds
for $14.4$ per cent of draws at $\sigma=0.8$ against $1.7$ per cent at
$\sigma=0.4$, which changes the shape of the distribution of $\theta_a$ and
raises $\E[\theta]$ from $1.00$ to $1.06$. We report both rows and draw no
conclusion from their difference.

\paragraph{Nothing here bears on the theorems.}
\paperref{thm:deriv}{Theorems}{2 and 3} of the paper hold for every network satisfying \paperref{ass:norm}{Assumption}{1}.
The table describes the measure of networks on which the two limits point in
opposite directions under one particular sampling scheme, which is a question of
prevalence and not of validity.

%=====================================================================
\section*{Replication}

The table is produced by \texttt{scripts/sign\_reversal.py} and the block
construction of \paperref{rem:blocks}{Remark}{5} by \texttt{scripts/block\_reversal.py}, in the
repository
%
\begin{center}
\url{https://github.com/renatovicente/price-of-risk-aggregation}
\end{center}
%
Both are self-contained, depend only on \texttt{numpy} and \texttt{scipy}, read
no data and draw from fixed seeds; \texttt{./reproduce.sh} runs them together
with the rest in about fifteen seconds.
